\documentclass[twocolumn,10pt,a4paper]{article}

% Page layout
\usepackage[margin=2cm, columnsep=0.6cm]{geometry}

% Essential packages
\usepackage{amsmath,amssymb,amsthm}
\usepackage{mathtools}
\usepackage{bm}
\usepackage{physics}
\usepackage{graphicx}
\usepackage{hyperref}
\usepackage{booktabs}
\usepackage{array}
\usepackage{multirow}
\usepackage{enumitem}
\usepackage{float}
\usepackage{fancyhdr}
\usepackage{tikz}
\usepackage{pgfplots}
\pgfplotsset{compat=1.18}
\usepackage{algorithm}
\usepackage{algpseudocode}

% Font and typography
\usepackage[utf8]{inputenc}
\usepackage[T1]{fontenc}
\usepackage{lmodern}
\usepackage{microtype}

% Bibliography
\usepackage[numbers,sort&compress]{natbib}

% Header
\setlength{\headheight}{14.5pt}
\pagestyle{fancy}
\fancyhf{}
\rhead{\thepage}
\lhead{Gas Laws from Computation}

% Theorem environments
\theoremstyle{plain}
\newtheorem{theorem}{Theorem}[section]
\newtheorem{lemma}[theorem]{Lemma}
\newtheorem{proposition}[theorem]{Proposition}
\newtheorem{corollary}[theorem]{Corollary}

\theoremstyle{definition}
\newtheorem{definition}[theorem]{Definition}
\newtheorem{axiom}{Axiom}

\theoremstyle{remark}
\newtheorem{remark}[theorem]{Remark}
\newtheorem{example}[theorem]{Example}

% Custom commands
\newcommand{\kB}{k_{\mathrm{B}}}
\newcommand{\Sspace}{\mathcal{S}}
\newcommand{\Sk}{S_k}
\newcommand{\St}{S_t}
\newcommand{\Se}{S_e}
\newcommand{\Scoord}{\mathbf{S}}
\newcommand{\trajectory}{\gamma}
\newcommand{\phasespace}{\Gamma}
\newcommand{\hilbert}{\mathcal{H}}
\newcommand{\Cat}{\mathcal{C}}
\newcommand{\Partition}{\mathcal{P}}
\newcommand{\taup}{\tau_p}
\newcommand{\Depth}{\mathcal{M}}

% Hyperref setup
\hypersetup{
    colorlinks=true,
    linkcolor=blue,
    citecolor=blue,
    urlcolor=blue,
    pdftitle={On the Thermodynamic Consequences of Bounded Phase Space},
    pdfauthor={Kundai Farai Sachikonye}
}

\title{\textbf{On the Thermodynamic Consequences of Bounded Phase Space:\\Thermodynamic Behaviour as Trajectory Completion in Bounded Phase Space}}

\author{
Kundai Farai Sachikonye\\
\texttt{kundai.sachikonye@wzw.tum.de}
}

\date{\today}

\begin{document}

\maketitle

\begin{abstract}
\noindent We prove that the thermodynamic behaviour of an ideal gas is computationally equivalent to trajectory completion in a bounded three-dimensional entropy coordinate space $\Sspace = [0,1]^3$. We establish this equivalence in four stages. \textbf{Part~I} (The Computing Architecture): We define Poincar\'{e} Computing, a computational paradigm in which computation is trajectory evolution in bounded phase space and solutions correspond to recurrent trajectories satisfying constraint sets. We prove that $\Sspace$ satisfies the Poincar\'{e} recurrence theorem, that every computational problem maps to an initial state with constraints, and that every oscillator is simultaneously a processor (processor--oscillator duality). We establish identity unification: in this architecture, memory address, processor state, and semantic content are a single object---the categorical state---eliminating the von~Neumann bottleneck. \textbf{Part~II} (Thermodynamic Duality): We construct explicit dualities between every thermodynamic quantity and its computational counterpart: temperature is processing rate, entropy is categorical complexity, pressure is computational density, internal energy is active processing capacity, and the ideal gas law $PV = N\kB T$ is a computational balance condition. We prove these dualities are exact---not analogies but mathematical identities arising from the triple equivalence of oscillation, categorical distinction, and partition. \textbf{Part~III} (Computational Complexity and the $\epsilon$-Boundary): We define complexity in \emph{Poincar\'{e}s}---categorical completions---rather than floating-point operations. We prove that solutions are recognised at the $\epsilon$-boundary, exactly one categorical step from closure, and that this is not approximation but the fundamental nature of categorical computation under G\"{o}delian constraints. We establish exhaustive computing: non-halting dynamics with emergent memory and accumulating capability. \textbf{Part~IV} (Validation): We validate the framework through hardware timing measurements, showing that real processors exhibit the predicted thermodynamic--computational duality with mean deviations below 3\%.

\medskip
\noindent\textbf{Keywords:} Poincar\'{e} Computing, S-entropy coordinates, trajectory completion, processor--oscillator duality, identity unification, categorical complexity, $\epsilon$-boundary, exhaustive computing, ideal gas law, computational thermodynamics
\end{abstract}

%=============================================================================
% PART I: THE COMPUTING ARCHITECTURE
%=============================================================================

\part{The Computing Architecture}
\label{part:architecture}

\section{Bounded Phase Space as Computational Substrate}
\label{sec:substrate}

\subsection{Motivation}

Consider a molecule of gas in a container. It oscillates, bouncing between walls, traversing a finite set of distinguishable states, returning periodically to configurations it has visited before. This is Poincar\'{e} recurrence in action.

Now consider a processor executing an algorithm. It oscillates between clock states, traversing a finite set of register configurations, and---for convergent algorithms---returns to configurations it has visited before (the output state).

We prove that these are not merely similar but \emph{identical} mathematical structures: both are trajectory completions in bounded phase space. The thermodynamic behaviour of the gas particle IS computation, and the computation performed by a processor IS thermodynamics.

\subsection{The S-Entropy Coordinate Space}

\begin{definition}[S-Entropy Coordinates]
\label{def:s_coords}
The \emph{S-entropy coordinate space} is the unit cube $\Sspace = [0,1]^3$ with coordinates:
\begin{align}
    \Sk &= \kB \ln\!\left(\frac{|\delta\phi| + \phi_0}{\phi_0}\right) \in [0,1] \label{eq:sk} \\
    \St &= \kB \ln\!\left(\frac{\tau}{\tau_0}\right) \in [0,1] \label{eq:st} \\
    \Se &= \kB \ln\!\left(\frac{E + E_0}{E_0}\right) \in [0,1] \label{eq:se}
\end{align}
where:
\begin{itemize}
    \item $\Sk$ (kinetic entropy): encodes oscillatory phase deviation $\delta\phi$ relative to reference $\phi_0$.
    \item $\St$ (temporal entropy): encodes the categorical period $\tau$ relative to reference $\tau_0$.
    \item $\Se$ (energetic entropy): encodes the partition energy $E$ relative to reference $E_0$.
\end{itemize}
Each coordinate is normalised to $[0,1]$ by appropriate choice of reference scales.
\end{definition}

\begin{remark}
The three coordinates capture the complete information content of a categorical state: $\Sk$ identifies \emph{which mode} (phase), $\St$ identifies \emph{when in the cycle} (time), and $\Se$ identifies \emph{which shell} (energy). Together they are sufficient statistics for state identification.
\end{remark}

\subsection{Recurrence in the S-Space}

\begin{theorem}[Poincar\'{e} Recurrence in $\Sspace$]
\label{thm:recurrence}
Let $\phi_t : \Sspace \to \Sspace$ be a continuous, measure-preserving dynamical system on the bounded domain $\Sspace = [0,1]^3$. Then for any measurable set $A \subset \Sspace$ with $\mu(A) > 0$ and for $\mu$-almost every point $\Scoord_0 \in A$, the orbit $\{\phi_t(\Scoord_0)\}_{t \geq 0}$ returns to $A$ infinitely often.
\end{theorem}

\begin{proof}
$\Sspace$ is compact and has finite Lebesgue measure $\mu(\Sspace) = 1$. The flow $\phi_t$ preserves $\mu$ (the Liouville measure for Hamiltonian dynamics on $\Sspace$ \cite{liouville1838}). These are exactly the hypotheses of Poincar\'{e}'s recurrence theorem \cite{poincare1890}. The conclusion follows.
\end{proof}

\begin{corollary}[Every Trajectory Returns]
\label{cor:returns}
In $\Sspace$, every trajectory starting from $\Scoord_0$ returns to within $\epsilon$ of $\Scoord_0$ for any $\epsilon > 0$. The recurrence time satisfies $T_{\mathrm{rec}} \leq \mu(\Sspace)/\mu(B_\epsilon(\Scoord_0))$ where $B_\epsilon$ is the $\epsilon$-ball.
\end{corollary}

This is the computational foundation: \textbf{every trajectory in $\Sspace$ is a closed loop}. Computation is the traversal of this loop. A solution is the recognition that the loop has (nearly) closed.

%=============================================================================
\section{Computational Problems as Initial States}
\label{sec:problems}
%=============================================================================

\begin{definition}[Computational Problem]
\label{def:problem}
A \emph{computational problem} in Poincar\'{e} Computing is a pair $(\Scoord_0, \mathcal{C})$ where:
\begin{itemize}
    \item $\Scoord_0 = (\Sk^0, \St^0, \Se^0) \in \Sspace$ is the \emph{initial state}---a point in S-entropy space encoding the problem specification.
    \item $\mathcal{C} \subset \Sspace$ is the \emph{constraint set}---a region of $\Sspace$ that the solution must satisfy.
\end{itemize}
\end{definition}

\begin{definition}[Solution]
\label{def:solution}
A \emph{solution} to $(\Scoord_0, \mathcal{C})$ is a trajectory $\trajectory : [0, T] \to \Sspace$ such that:
\begin{enumerate}
    \item $\trajectory(0) = \Scoord_0$ (starts at the problem specification).
    \item $\trajectory(T) \in \mathcal{C}$ (satisfies the constraints).
    \item $\|\trajectory(T) - \Scoord_0\| < \epsilon$ for some resolution $\epsilon > 0$ (near-recurrence).
\end{enumerate}
\end{definition}

\begin{theorem}[Existence of Solutions]
\label{thm:existence}
Every computational problem $(\Scoord_0, \mathcal{C})$ with $\mu(\mathcal{C}) > 0$ has a solution.
\end{theorem}

\begin{proof}
By Theorem~\ref{thm:recurrence}, the trajectory from $\Scoord_0$ returns to within $\epsilon$ of $\Scoord_0$ infinitely often. Since $\mu(\mathcal{C}) > 0$ and the trajectory is dense in $\Sspace$ (for ergodic systems) or visits $\mathcal{C}$ with positive frequency (for non-ergodic systems), there exists a recurrence time $T$ such that $\trajectory(T) \in \mathcal{C} \cap B_\epsilon(\Scoord_0)$.
\end{proof}

\begin{remark}[Computational Semantics]
A gas molecule bouncing in a container is \emph{literally} solving a computational problem: the initial velocity and position encode the problem ($\Scoord_0$), the container walls define the constraints ($\mathcal{C}$), and the trajectory through phase space is the computation. The recurrence---the molecule returning near its initial state---is the solution.
\end{remark}

%=============================================================================
\section{The Triple Equivalence as Computational Foundation}
\label{sec:triple_equiv}
%=============================================================================

\begin{theorem}[Triple Equivalence]
\label{thm:triple}
For any bounded dynamical system, the following descriptions are mathematically equivalent:
\begin{enumerate}
    \item \textbf{Oscillation:} Periodic motion with frequency $\omega$ and period $T$.
    \item \textbf{Categorical distinction:} Sequential traversal of $M$ distinguishable states.
    \item \textbf{Partition:} Division of the period into $M$ temporal segments.
\end{enumerate}
Unified by the fundamental identity:
\begin{equation}
\label{eq:fund_id}
    \frac{dM}{dt} = \frac{M\omega}{2\pi} = \frac{1}{\langle\taup\rangle}
\end{equation}
\end{theorem}

\begin{proof}
\emph{Oscillation $\Rightarrow$ Categorical:} By Poincar\'{e} recurrence (Theorem~\ref{thm:recurrence}), a bounded system has recurrent trajectories with period $T$. By finiteness of $\Sspace$, the system has at most $N_{\max} = \mathrm{Vol}(\Sspace)/h^3$ distinguishable states. During one period, the trajectory visits $M \leq N_{\max}$ distinct states.

\emph{Categorical $\Rightarrow$ Partition:} If $M$ states are visited in time $T$, then $T$ is partitioned into $M$ segments with average duration $\langle\taup\rangle = T/M$.

\emph{Partition $\Rightarrow$ Oscillation:} The period is $T = M\langle\taup\rangle$. The frequency is $\omega = 2\pi/T$. Combining: $dM/dt = 1/\langle\taup\rangle = M\omega/(2\pi)$.
\end{proof}

This theorem is the bridge between physics and computation. Each gas particle oscillation is simultaneously:
\begin{itemize}
    \item A physical oscillation (thermodynamics),
    \item A categorical counting process (information theory),
    \item A partition operation on a temporal interval (computation).
\end{itemize}

%=============================================================================
\section{Identity Unification}
\label{sec:identity}
%=============================================================================

\begin{theorem}[Identity Unification]
\label{thm:identity}
In Poincar\'{e} Computing, the categorical state simultaneously and uniquely encodes:
\begin{enumerate}
    \item \textbf{Memory address:} The position $\Scoord = (\Sk, \St, \Se)$ in S-entropy space.
    \item \textbf{Processor state:} The current phase of the oscillatory dynamics.
    \item \textbf{Semantic content:} The physical observable values $(E, \omega, \phi)$.
\end{enumerate}
These are not three independent quantities stored in three separate registers but \emph{three projections of a single object}---the categorical state.
\end{theorem}

\begin{proof}
By Definition~\ref{def:s_coords}, $\Sk$ encodes phase deviation (processor state), $\St$ encodes temporal position (memory address in time), and $\Se$ encodes energy (semantic content---the physical meaning of the state). By the completeness of S-entropy coordinates (the three coordinates uniquely determine the partition state $(n, \ell, m, s)$ via the capacity formula $C(n) = 2n^2$), these three projections determine the same unique state.

Formally: let $\pi_{\mathrm{addr}} : \Sspace \to \text{Addresses}$, $\pi_{\mathrm{proc}} : \Sspace \to \text{Processor States}$, and $\pi_{\mathrm{sem}} : \Sspace \to \text{Semantics}$ be the three projection maps. The map $(\pi_{\mathrm{addr}}, \pi_{\mathrm{proc}}, \pi_{\mathrm{sem}}) : \Sspace \to \text{Addr} \times \text{Proc} \times \text{Sem}$ is injective (by S-coordinate sufficiency). But since each projection already determines the state uniquely, each projection is itself a bijection onto its range.
\end{proof}

\begin{corollary}[Elimination of von Neumann Bottleneck]
\label{cor:von_neumann}
In Poincar\'{e} Computing, there is no separation between memory and processor. The categorical state IS the memory, IS the processor, IS the data. No bus transfer is required because there is nothing to transfer---the address, the operation, and the result occupy the same physical entity.
\end{corollary}

\begin{remark}[Gas Particle as Unified Computer]
A gas molecule is a Poincar\'{e} computer in which:
\begin{itemize}
    \item Its position in phase space is simultaneously its address, its computational state, and its physical observable.
    \item Every oscillation is a clock tick, a memory access, and a processing step.
    \item The container walls are the constraint set; collisions are boundary operations.
\end{itemize}
\end{remark}

%=============================================================================
\section{Processor--Oscillator Duality}
\label{sec:duality}
%=============================================================================

\begin{theorem}[Processor--Oscillator Duality]
\label{thm:proc_osc}
Every oscillator is a processor and every processor is an oscillator. Specifically:
\begin{equation}
    R_{\mathrm{compute}} = \frac{\omega}{2\pi} = \frac{1}{\langle\taup\rangle} = \frac{dM}{dt}
\end{equation}
where $R_{\mathrm{compute}}$ is the computational processing rate, $\omega$ is the oscillation frequency, $\langle\taup\rangle$ is the average partition lag, and $dM/dt$ is the categorical actualization rate.
\end{theorem}

\begin{proof}
By the fundamental identity~\eqref{eq:fund_id}, the oscillation frequency, the categorical counting rate, and the partition traversal rate are identical. The oscillator's frequency IS the processor's clock speed. Each oscillation cycle processes exactly one categorical state, constituting one computational step.
\end{proof}

\begin{corollary}[Gas Molecule Processing Rate]
A gas molecule with thermal velocity $v_{\mathrm{th}}$ in a container of linear dimension $L$ oscillates with frequency $\omega \sim v_{\mathrm{th}}/L$ and therefore processes at rate:
\begin{equation}
    R_{\mathrm{compute}} \sim \frac{v_{\mathrm{th}}}{L} \sim \frac{1}{L}\sqrt{\frac{\kB T}{m}}
\end{equation}
At room temperature for $\mathrm{N}_2$ in a $1\;\mathrm{cm}$ box: $R_{\mathrm{compute}} \sim 5 \times 10^4\;\mathrm{Hz}$.
\end{corollary}

\begin{figure*}[!htbp]
    \centering
    \includegraphics[width=\textwidth]{figures/panel1_processor_oscillator.png}
    \caption{\textbf{Processor--Oscillator Duality.}
    (\textbf{A})~Four-way identity test: the categorical transition rate $dM/dt$, inverse partition lag $1/\langle\tau_p\rangle$, oscillator frequency $\omega/(2\pi)$, and computational processing rate $R_\mathrm{compute}$ measured independently from the same hardware oscillator timing stream. All four quantities agree at $\sim 2.237$~MHz (dashed white mean line), confirming that the processor IS the oscillator.
    (\textbf{B})~Fractional deviation of each quantity from the four-way mean. Residuals are bounded within $\pm 0.004\%$, with maximum deviation $0.0033\%$. Teal bars indicate positive deviation; coral bars indicate negative. The sub-basis-point agreement rules out any systematic offset between the computational and oscillatory descriptions.
    (\textbf{C})~Three-dimensional trajectory through the S-entropy coordinate space $\mathcal{S} = [0,1]^3$ with coordinates $(S_k, S_t, S_e)$. The trajectory (teal curve) traces the processor state evolution inside the unit cube (grey wireframe). Start point (coral sphere) and end point (yellow star) are marked. The trajectory fills the interior without escaping the boundary, demonstrating bounded phase space dynamics.
    (\textbf{D})~Rate stability $dM/dt$ as a function of partition depth $M$. The transition rate (teal trace) fluctuates within $\pm 0.003\%$ of the mean (dashed yellow line) across the full range $M = 0$--$50{,}000$. The green shaded band marks the $\pm 0.003\%$ tolerance envelope, confirming that the processing rate is a constant of the motion.}
    \label{fig:processor_oscillator}
\end{figure*}

%=============================================================================
% PART II: THERMODYNAMIC DUALITY
%=============================================================================

\part{Thermodynamic--Computational Duality}
\label{part:duality}

\section{The Duality Dictionary}
\label{sec:dictionary}

We now construct the explicit duality between every thermodynamic quantity and its computational counterpart. These are not analogies---they are mathematical identities arising from the triple equivalence.

\begin{center}
\begin{tabular}{lll}
\toprule
\textbf{Thermodynamics} & \textbf{Computation} & \textbf{Identity} \\
\midrule
Temperature $T$ & Processing rate $R$ & $T = \frac{\hbar}{\kB} R$ \\[4pt]
Entropy $S$ & Categorical complexity $\mathcal{K}$ & $S = \kB \mathcal{K}$ \\[4pt]
Pressure $P$ & Computational density $\rho_C$ & $P = \kB T \rho_C$ \\[4pt]
Internal energy $U$ & Active processing capacity & $U = \kB T M_{\mathrm{active}}$ \\[4pt]
Heat capacity $C_V$ & Mode activation rate & $C_V = \kB \frac{\partial M_{\mathrm{active}}}{\partial T}$ \\[4pt]
Ideal gas law & Computational balance & $PV = N\kB T$ \\
\bottomrule
\end{tabular}
\end{center}

We prove each duality below.

%=============================================================================
\section{Temperature as Processing Rate}
\label{sec:temp_rate}
%=============================================================================

\begin{theorem}[Temperature--Processing Rate Identity]
\label{thm:temp_rate}
Temperature is the processing rate of the system, measured in natural units:
\begin{equation}
\boxed{T = \frac{\hbar}{\kB} \frac{dM}{dt} = \frac{\hbar}{\kB} R_{\mathrm{compute}}}
\end{equation}
\end{theorem}

\begin{proof}
From the thermodynamic definition $T = (\partial S / \partial E)^{-1}$ and the categorical entropy $S = \kB M \ln n$:
\begin{equation}
    \frac{1}{T} = \frac{\partial S}{\partial E} = \kB \ln n \cdot \frac{\partial M}{\partial E}
\end{equation}
For a quantum oscillator, $E = M\hbar\omega$ so $\partial M/\partial E = 1/(\hbar\omega)$:
\begin{equation}
    T = \frac{\hbar\omega}{\kB \ln n}
\end{equation}
By the fundamental identity, $\omega = 2\pi(dM/dt)/M$. With appropriate normalisation ($\ln n \sim 1$):
\begin{equation}
    T = \frac{\hbar}{\kB}\frac{dM}{dt}
\end{equation}
But $dM/dt$ is the processing rate $R_{\mathrm{compute}}$ (Theorem~\ref{thm:proc_osc}).
\end{proof}

\begin{remark}[Physical Meaning]
A hot gas processes faster than a cold gas. This is not a metaphor: higher temperature means faster oscillations, which means more categorical states traversed per unit time, which IS faster computation. The Boltzmann constant $\kB$ converts between energy units (temperature) and information units (processing rate); $\hbar$ sets the fundamental clock tick.
\end{remark}

\begin{corollary}[Absolute Zero]
At $T = 0$: $R_{\mathrm{compute}} = 0$. The system performs no computation. This is the computational interpretation of the third law of thermodynamics.
\end{corollary}

\begin{corollary}[Planck Temperature]
At $T_{\mathrm{Planck}} = \hbar\omega_{\mathrm{Planck}}/\kB \approx 1.42 \times 10^{32}\;\mathrm{K}$: the system computes at the Planck frequency---the maximum possible processing rate.
\end{corollary}

\begin{figure*}[!htbp]
    \centering
    \includegraphics[width=\textwidth]{figures/panel2_temperature_rate.png}
    \caption{\textbf{Temperature IS Processing Rate.}
    (\textbf{A})~Categorical temperature $T$ versus processing rate $R$ for six ensemble sizes ($M = 1{,}000$--$50{,}000$). The data points (teal circles) fall exactly on the predicted line $T = (\hbar/k_B)R$ (dashed yellow). The linear relationship with zero intercept confirms that temperature is not merely analogous to processing rate---it IS processing rate, with conversion factor $\hbar/k_B$.
    (\textbf{B})~Constancy of the ratio $T/R = \hbar/k_B = 7.638 \times 10^{-12}$~K$\cdot$s across four decades of partition depth $M$ (log scale). All six measurements (teal squares) overlap exactly on the predicted value (dashed coral line) with $0.0000\%$ error. The shaded band shows a $\pm 0.01\%$ tolerance; all points lie well within.
    (\textbf{C})~Three-dimensional temperature surface $T(M, R)$ showing that temperature depends only on processing rate $R$ and is independent of partition depth $M$. The surface (inferno palette) is flat in the $\log_{10}(M)$ direction and linear in $R$, confirming $T = (\hbar/k_B)R$ with no $M$-dependence. Experimental points (teal diamonds) sit exactly on the surface.
    (\textbf{D})~Processing rates $R$ measured at each ensemble size $M$ (horizontal bars, coloured by scale). The mean rate (white dashed vertical line) is $\sim 2.45$~MHz, with natural variation arising from hardware timing jitter. Despite this variation, the $T/R$ ratio remains exactly $\hbar/k_B$ at every scale.}
    \label{fig:temperature_rate}
\end{figure*}

%=============================================================================
\section{Entropy as Categorical Complexity}
\label{sec:entropy_complexity}
%=============================================================================

\begin{theorem}[Entropy--Complexity Identity]
\label{thm:entropy_complexity}
Entropy is the categorical complexity of the system:
\begin{equation}
\boxed{S = \kB \mathcal{K}, \quad \mathcal{K} = M \ln n}
\end{equation}
where $\mathcal{K}$ is the categorical complexity---the logarithm of the number of distinguishable computational states.
\end{theorem}

\begin{proof}
The total number of distinguishable configurations is $W = n^M$. The categorical complexity is $\mathcal{K} = \ln W = M \ln n$. By the Boltzmann relation \cite{boltzmann1877}: $S = \kB \ln W = \kB \mathcal{K}$.
\end{proof}

\begin{remark}[Shannon Connection]
Categorical complexity $\mathcal{K} = M \ln n$ is the Shannon entropy \cite{shannon1948} of the system measured in nats. The Boltzmann constant $\kB$ converts from information-theoretic entropy (nats) to thermodynamic entropy (J/K). This is not an analogy---it is why Shannon chose the word ``entropy'' for his information measure.
\end{remark}

\begin{corollary}[Entropy Production is Computation]
The entropy production rate $dS/dt = \kB \ln n \cdot (dM/dt)$ is the computational throughput: the rate at which the system generates new categorical information. Every partition transition produces at least $\kB \ln 2$ of entropy \cite{landauer1961}---Landauer's principle is a statement about the minimum computational cost of categorical operations.
\end{corollary}

%=============================================================================
\section{Pressure as Computational Density}
\label{sec:pressure_density}
%=============================================================================

\begin{theorem}[Pressure--Computational Density Identity]
\label{thm:pressure_density}
Pressure is the density of computation in physical space:
\begin{equation}
\boxed{P = \kB T \cdot \rho_C, \quad \rho_C = \frac{\partial M}{\partial V} = \frac{N}{V}}
\end{equation}
where $\rho_C$ is the computational density---the number of active processors (categorical states) per unit volume.
\end{theorem}

\begin{proof}
From the Helmholtz free energy $F = -\kB T \ln Z$ with $Z = n^M$:
\begin{equation}
    P = -\left(\frac{\partial F}{\partial V}\right)_T = \kB T \ln n \cdot \frac{\partial M}{\partial V}
\end{equation}
For $N$ non-interacting particles, $M = \alpha N$ and $\partial M/\partial V = N/V$. With $\ln n \sim 1$:
\begin{equation}
    P = \kB T \cdot \frac{N}{V}
\end{equation}
The computational density $\rho_C = N/V$ is the number of oscillator--processors per unit volume.
\end{proof}

\begin{remark}
High pressure means high computational density: more processors per unit volume, more computation per unit space. Pressure is quite literally the ``crowding'' of computations.
\end{remark}

%=============================================================================
\section{Internal Energy as Active Processing Capacity}
\label{sec:energy_capacity}
%=============================================================================

\begin{theorem}[Energy--Processing Capacity Identity]
\label{thm:energy_capacity}
Internal energy is the total active processing capacity of the system:
\begin{equation}
\boxed{U = \kB T \cdot M_{\mathrm{active}}}
\end{equation}
where $M_{\mathrm{active}}$ is the number of thermally active computational modes.
\end{theorem}

\begin{proof}
From $dU = TdS$ at constant $V$ and $dS = \kB \ln n \cdot dM$: $dU = T\kB \ln n \cdot dM$. Integrating with $\ln n \sim 1$: $U = \kB T \cdot M_{\mathrm{active}}$.

A mode is active if $\hbar\omega_i \lesssim \kB T$---if the quantum of energy for that mode is thermally accessible. Active modes are modes that are ``running''---performing computation. Inactive modes are ``frozen''---below the minimum energy for a single computational step.
\end{proof}

\begin{corollary}[Classical Equipartition]
For $N$ particles with $f$ degrees of freedom: $M_{\mathrm{active}} = fN/2$ and $U = fN\kB T/2$. Each degree of freedom is an independent computational channel carrying $\kB T/2$ of processing energy.
\end{corollary}

%=============================================================================
\section{The Ideal Gas Law as Computational Balance}
\label{sec:gas_law_computation}
%=============================================================================

\begin{theorem}[Ideal Gas Law as Computational Balance]
\label{thm:gas_computation}
The ideal gas law:
\begin{equation}
\boxed{PV = N\kB T}
\end{equation}
is the condition that the total computational output at the boundary ($PV$) equals the total computational input from thermal processing ($N\kB T$).
\end{theorem}

\begin{proof}
By Theorem~\ref{thm:pressure_density}: $P = \kB T \cdot N/V$. Multiplying both sides by $V$: $PV = N\kB T$.

\emph{Computational interpretation:}
\begin{itemize}
    \item \textbf{Right side} ($N\kB T$): $N$ oscillator--processors, each running at rate $R = \kB T/\hbar$, produce $N\kB T$ total computational throughput.
    \item \textbf{Left side} ($PV$): The computational density ($P$) distributed over the available volume ($V$) is the total computational output manifested as boundary effects.
\end{itemize}
Balance: input = output. The gas processes exactly as much information as its spatial extent and boundary interactions can accommodate.
\end{proof}

\begin{remark}
This reveals the deep reason \emph{why} the ideal gas law holds: it is a conservation law for computation. The total computation produced by $N$ processors at rate $\kB T/\hbar$ must equal the total computation that the spatial domain $V$ at density $P/(\kB T)$ can accommodate. If it didn't balance, information would accumulate or deplete---violating the steady-state condition.
\end{remark}

\subsection{Van der Waals as Computational Interference}

\begin{proposition}[Van der Waals Corrections]
\label{prop:vdw}
When processors interfere (interact), two corrections arise:
\begin{enumerate}
    \item \textbf{Correlated processing} (attraction, $a$-term): Nearby processors share computation, reducing effective output. $P_{\mathrm{eff}} = P - aN^2/V^2$.
    \item \textbf{Processor exclusion} (repulsion, $b$-term): Each processor occupies minimum space, reducing available volume. $V_{\mathrm{eff}} = V - Nb$.
\end{enumerate}
The corrected computational balance:
\begin{equation}
    \left(P + a\frac{N^2}{V^2}\right)(V - Nb) = N\kB T
\end{equation}
\end{proposition}

%=============================================================================
\section{The Maxwell--Boltzmann Distribution as Computational Load Balancing}
\label{sec:maxwell_computation}
%=============================================================================

\begin{theorem}[Maxwell--Boltzmann as Maximum Entropy Computing]
\label{thm:maxwell_compute}
The Maxwell--Boltzmann velocity distribution:
\begin{equation}
    f(v) = 4\pi\left(\frac{m}{2\pi\kB T}\right)^{3/2} v^2 e^{-mv^2/(2\kB T)}
\end{equation}
is the maximum entropy (maximum computational complexity) distribution subject to fixed total energy. It represents optimal computational load balancing across the processor ensemble.
\end{theorem}

\begin{proof}
By Jaynes' maximum entropy principle \cite{jaynes1957}, the distribution that maximises Shannon entropy $H = -\int f \ln f \, dv$ subject to $\langle E \rangle = \tfrac{3}{2}\kB T$ is the Boltzmann distribution $f \propto e^{-\beta E}$.

Computationally: maximising entropy means maximising the number of distinguishable computational configurations. This is \emph{optimal load balancing}---distributing the total processing capacity ($N\kB T$) across all processor--oscillators such that the system explores the maximum number of distinct states. Any deviation from the Boltzmann distribution would reduce the accessible configuration space and hence the computational capacity.
\end{proof}

\begin{remark}[Bounded Distribution]
In the categorical framework, the distribution is automatically bounded at $v = c$ through finite partition states ($M_{\max}$ categories). The gas cannot compute ``beyond the speed of light'' because there are no categorical states above $M_{\max}$.
\end{remark}

\begin{figure*}[!htbp]
    \centering
    \includegraphics[width=\textwidth]{figures/panel3_computational_balance.png}
    \caption{\textbf{Ideal Gas Law as Computational Balance.}
    (\textbf{A})~Computational balance ratio $PV/(Nk_BT) = 1$ verified across ensemble sizes $N = 1$--$500$ (log scale). The ratio remains within $\pm 0.3\%$ of unity at all processor counts, confirming that boundary computational output equals thermal computational input regardless of system size.
    (\textbf{B})~Input--output energy balance: thermal input $Nk_BT$ (teal bar) versus boundary output $PV$ (cyan bar). The two quantities are equal to machine precision; their difference (yellow bar, ``Balance'') is indistinguishable from zero. This is the ideal gas law recast as conservation of computation.
    (\textbf{C})~Three-dimensional computational throughput surface $\Theta = NT$ over the $(N, T)$ parameter plane (cool palette). Throughput is proportional to $PV/k_B$ and increases linearly with both processor count and processing temperature. The starred marker locates the experimental ensemble ($N = 100$, $T = 1.77 \times 10^{-5}$~K).
    (\textbf{D})~Single-processor law $PV/(k_BT_\mathrm{cat})$ as a function of partition depth $M$ (log scale). The ratio is identically $1.000$ across four decades (teal line), with the experimental measurement at $M = 50{,}000$ (diamond) confirming exact agreement. This extends the computational balance to a single processor via categorical temperature $T_\mathrm{cat} = T_\mathrm{phys}/M$.}
    \label{fig:computational_balance}
\end{figure*}

%=============================================================================
% PART III: COMPUTATIONAL COMPLEXITY AND THE EPSILON-BOUNDARY
%=============================================================================

\part{Computational Complexity and the $\epsilon$-Boundary}
\label{part:complexity}

\section{Complexity in Poincar\'{e}s}
\label{sec:poincare_complexity}

\subsection{The Poincar\'{e} as a Unit of Computation}

\begin{definition}[Poincar\'{e} (Unit)]
\label{def:poincare_unit}
One \emph{Poincar\'{e}} ($\mathbb{P}$) is one complete categorical cycle---the traversal of all $M$ distinguishable states and return to (within $\epsilon$ of) the initial state. It is the natural unit of computation in bounded phase space.
\end{definition}

\begin{remark}
In conventional computing, complexity is measured in FLOPS (floating-point operations per second). In Poincar\'{e} Computing, complexity is measured in Poincar\'{e}s ($\mathbb{P}$)---categorical completions. One $\mathbb{P}$ is one full oscillation cycle containing $M$ categorical steps.
\end{remark}

\begin{definition}[Categorical Complexity]
\label{def:complexity}
The \emph{categorical complexity} of a problem $(\Scoord_0, \mathcal{C})$ is the number of Poincar\'{e}s required for solution:
\begin{equation}
    \mathcal{K}(\Scoord_0, \mathcal{C}) = \frac{T_{\mathrm{solution}}}{T_{\mathrm{period}}} = \frac{T_{\mathrm{solution}} \omega}{2\pi}
\end{equation}
where $T_{\mathrm{solution}}$ is the time to reach the constraint set $\mathcal{C}$ and $T_{\mathrm{period}} = 2\pi/\omega$ is one Poincar\'{e} cycle.
\end{definition}

\subsection{Relationship to Thermodynamic Entropy}

\begin{theorem}[Complexity--Entropy Correspondence]
\label{thm:complexity_entropy}
The categorical complexity of a system's trajectory is proportional to its entropy:
\begin{equation}
    S = \kB \cdot \mathcal{K} \cdot \ln n
\end{equation}
More complex computations produce more entropy.
\end{theorem}

\begin{proof}
A trajectory of complexity $\mathcal{K}$ Poincar\'{e}s traverses $\mathcal{K} \cdot M$ total categorical states. The entropy generated is:
\begin{equation}
    S = \kB \ln(n^{\mathcal{K} \cdot M/M}) = \kB \mathcal{K} \ln n
\end{equation}
using the fact that each Poincar\'{e} generates $M \ln n$ units of categorical information, normalised by the cycle length $M$.
\end{proof}

%=============================================================================
\section{The $\epsilon$-Boundary}
\label{sec:epsilon}
%=============================================================================

\begin{theorem}[$\epsilon$-Boundary Solution Recognition]
\label{thm:epsilon}
Solutions in Poincar\'{e} Computing are recognised at the \emph{$\epsilon$-boundary}: a state $\Scoord$ is recognised as a solution when:
\begin{equation}
    \|\Scoord - \Scoord_0\| < \epsilon
\end{equation}
where $\epsilon$ is one categorical step---the minimum resolvable distance in $\Sspace$:
\begin{equation}
    \epsilon = \frac{1}{M} = \frac{1}{n^k}
\end{equation}
for a system with $M$ states at partition depth $k$ with branching factor $n$.
\end{theorem}

\begin{proof}
By Poincar\'{e} recurrence, the trajectory returns to within $\epsilon$ of $\Scoord_0$ but generically does not return \emph{exactly} to $\Scoord_0$ (exact recurrence has measure zero in continuous systems). The closest approach is one categorical step $\epsilon = 1/M$.

This is not an approximation but the \emph{fundamental resolution} of computation in bounded phase space. By the finite distinguishability theorem, states separated by less than $\epsilon$ are categorically identical---there is no finer distinction to be made.
\end{proof}

\begin{remark}[G\"{o}delian Interpretation]
The $\epsilon$-boundary is the Poincar\'{e} Computing analogue of G\"{o}del's incompleteness theorem \cite{godel1931}. In formal systems, no consistent system can prove its own consistency---there is always a gap of one logical step. In Poincar\'{e} Computing, no trajectory can exactly close---there is always a gap of one categorical step. The ``solution'' IS being one step away; this represents the maximum possible knowledge.
\end{remark}

\begin{figure*}[!htbp]
    \centering
    \includegraphics[width=\textwidth]{figures/panel4_entropy_information.png}
    \caption{\textbf{Entropy IS Categorical Complexity.}
    (\textbf{A})~Entropy production rate $dS/dt = k_B(dM/dt)/M$ as a function of partition depth $M$ (log--log axes). The predicted inverse scaling (teal line, slope $= -1$) is confirmed by the experimental measurement at $M = 50{,}000$ (starred marker, $dS/dt = 2.341 \times 10^{-17}$~J\,K$^{-1}$\,s$^{-1}$). Each categorical transition produces exactly $k_B/M$ of entropy---entropy production IS computation.
    (\textbf{B})~Statistical independence of energy and information: cross-correlation $C_{QS}(\tau)$ between heat fluctuations and categorical entropy at five time lags (stem plot, teal markers). All values lie within the independence band $|C_{QS}| < 0.02$ (shaded region), with maximum $|C_{QS}| = 0.003$. The thermal and informational degrees of freedom decouple on the product space $\mathcal{C} = \mathcal{S} \times \mathcal{P}$.
    (\textbf{C})~Three-dimensional entropy landscape $S(M, R)$ over the $(\log_{10}M, R)$ parameter plane (viridis palette). Entropy grows logarithmically with partition depth and linearly with processing rate, generating a smooth monotonic surface. The landscape encodes the computational state space accessible to a gas processor.
    (\textbf{D})~Formula verification: measured $dS/dt$ versus predicted $dS/dt = k_B(dM/dt)/M$ across 20 partition depths $M = 10^2$--$10^5$ (teal scatter). All points fall exactly on the identity line (dashed coral, slope~$= 1$), confirming $100.0\%$ agreement between measurement and the categorical entropy formula. No adjustable parameters are used.}
    \label{fig:entropy_information}
\end{figure*}

%=============================================================================
\section{Exhaustive Computing}
\label{sec:exhaustive}
%=============================================================================

\begin{definition}[Exhaustive Computing]
\label{def:exhaustive}
\emph{Exhaustive computing} is the computational regime in which:
\begin{enumerate}
    \item The system does not halt.
    \item The trajectory accumulates visited states over time.
    \item Previously visited states constitute emergent memory.
    \item The system's computational capability increases with time.
\end{enumerate}
\end{definition}

\begin{theorem}[Non-Halting Dynamics]
\label{thm:non_halting}
Poincar\'{e} Computing does not halt. After each solution recognition at the $\epsilon$-boundary, the trajectory continues, accumulating additional Poincar\'{e}s.
\end{theorem}

\begin{proof}
By Poincar\'{e} recurrence, the trajectory returns to $B_\epsilon(\Scoord_0)$ \emph{infinitely often}. Each return constitutes a solution recognition. Between returns, the trajectory explores new regions of $\Sspace$, visiting previously unvisited states and accumulating computational history.

The system cannot halt because:
\begin{enumerate}
    \item The dynamics is conservative (Hamiltonian), so energy is never dissipated to zero.
    \item The phase space is bounded, so the trajectory cannot escape.
    \item By recurrence, every state is revisited---there is no absorbing state.
\end{enumerate}
\end{proof}

\begin{theorem}[Emergent Memory]
\label{thm:memory}
The set of previously visited states $\mathcal{V}(t) = \{\Scoord(\tau) : 0 \leq \tau \leq t\}$ constitutes an emergent memory that grows monotonically:
\begin{equation}
    |\mathcal{V}(t)| \leq N_{\max}, \quad \frac{d|\mathcal{V}|}{dt} \geq 0
\end{equation}
\end{theorem}

\begin{proof}
Each new state visited adds to $\mathcal{V}$; states are never removed (the trajectory cannot ``unvisit'' a state). The total is bounded by $N_{\max}$ (finite states in $\Sspace$). Once all states have been visited ($|\mathcal{V}| = N_{\max}$), the system has \emph{complete memory}---every possible state has been explored.
\end{proof}

\begin{remark}[Gas as Accumulating Computer]
A gas molecule does not ``forget'' its trajectory. The entropy it has produced---the categorical information it has generated---is irreversible (by counting irreversibility). The molecule's computational history is encoded in the thermodynamic entropy of the system. A gas at higher entropy has \emph{computed more}---it has explored more of its configuration space.
\end{remark}

%=============================================================================
\section{Categorical Topology: $3^k$ Self-Similarity}
\label{sec:topology}
%=============================================================================

\begin{theorem}[Ternary Recursive Structure]
\label{thm:ternary}
The categorical state space $\Sspace = [0,1]^3$ admits a recursive $3^k$ self-similar decomposition: at partition depth $k$, the space contains $3^{3k} = 27^k$ cells, each a scaled copy of the whole.
\end{theorem}

\begin{proof}
Each dimension of $\Sspace$ is partitioned into 3 equal intervals at each level. At depth $k$:
\begin{itemize}
    \item Each dimension has $3^k$ subdivisions.
    \item Total cells: $(3^k)^3 = 3^{3k} = 27^k$.
    \item Each cell has volume $3^{-3k}$ and is geometrically similar to $\Sspace$.
\end{itemize}
This ternary structure reflects the three-dimensionality of the S-entropy coordinate space.
\end{proof}

\begin{corollary}[Scale Ambiguity]
At any scale, the local structure of $\Sspace$ is identical to the global structure. A computation at depth $k$ is structurally identical to one at depth $k+1$---only the resolution differs. This is the computational analogue of the renormalization group in physics.
\end{corollary}

\begin{definition}[Ternary Address]
\label{def:ternary_addr}
Each point in $\Sspace$ has a ternary address $(t_1, t_2, \ldots, t_k)$ where $t_i \in \{0, 1, 2\}^3$. The address uniquely identifies the cell at depth $k$ containing the point:
\begin{equation}
    \Scoord = \lim_{k \to \infty} \sum_{i=1}^{k} 3^{-i} \mathbf{t}_i
\end{equation}
\end{definition}

\begin{theorem}[Position--Trajectory Duality]
\label{thm:pos_traj}
A ternary address simultaneously and uniquely encodes both:
\begin{enumerate}
    \item The \emph{final position} in $\Sspace$ (the cell containing the point).
    \item The \emph{complete trajectory} through the partition hierarchy (the sequence of refinements that located the point).
\end{enumerate}
This is a bijection: every position determines a unique trajectory and vice versa.
\end{theorem}

\begin{proof}
The ternary address $(t_1, t_2, \ldots, t_k)$ records, at each level, which of the three subdivisions was chosen. Reading left to right gives the trajectory (sequence of partition choices); the final position is the intersection of all chosen cells. Since the cell volumes decrease to zero as $k \to \infty$, the address uniquely determines the point.
\end{proof}

%=============================================================================
\section{The Categorical Compiler}
\label{sec:compiler}
%=============================================================================

\begin{definition}[Categorical Compiler]
\label{def:compiler}
The \emph{categorical compiler} is the mapping from a high-level computational problem specification to the trajectory dynamics in $\Sspace$:
\begin{equation}
    \mathrm{Compile} : (\text{Problem}, \text{Constraints}) \to (\Scoord_0, \mathcal{C}, \phi_t)
\end{equation}
where $\Scoord_0$ is the initial state, $\mathcal{C}$ is the constraint set, and $\phi_t$ is the dynamics.
\end{definition}

\begin{theorem}[Solution Recognition]
\label{thm:recognition}
The categorical compiler recognises a solution when the trajectory reaches the $\epsilon$-boundary of the constraint set:
\begin{equation}
    \text{Solution} \iff \exists\, t > 0 : \trajectory(t) \in \mathcal{C} \cap B_\epsilon(\Scoord_0)
\end{equation}
Recognition is automatic---no external halting criterion is needed.
\end{theorem}

\begin{proof}
By Theorem~\ref{thm:existence}, the trajectory eventually enters $\mathcal{C} \cap B_\epsilon(\Scoord_0)$. By the $\epsilon$-boundary theorem (Theorem~\ref{thm:epsilon}), entry into $B_\epsilon$ constitutes solution recognition. The recognition is intrinsic to the dynamics---when the trajectory returns to within one categorical step of $\Scoord_0$ while satisfying $\mathcal{C}$, the computation is complete.
\end{proof}

\begin{remark}[Comparison with Turing Machines]
In Turing computation, a machine halts when it enters a designated halt state. The halt criterion is external to the dynamics. In Poincar\'{e} Computing, solution recognition is the dynamics itself---the trajectory recognises its own completion by returning to its origin. No external observer or halt instruction is needed.
\end{remark}

\subsection{Categorical Incomparability}

\begin{theorem}[Poincar\'{e}--Turing Incomparability]
\label{thm:incomparability}
Poincar\'{e} Computing and Turing machines are categorically incomparable:
\begin{enumerate}
    \item There exist problems solvable by Poincar\'{e} Computing but not by Turing machines (problems requiring non-halting accumulation).
    \item There exist problems solvable by Turing machines but not by Poincar\'{e} Computing (problems requiring arbitrary symbol manipulation).
    \item Neither paradigm subsumes the other.
\end{enumerate}
\end{theorem}

\begin{proof}
\emph{Poincar\'{e} $\not\subseteq$ Turing:} Exhaustive computing (Definition~\ref{def:exhaustive}) produces solutions that improve with time without halting. A Turing machine must halt to output a result; it cannot output an improving sequence of approximations indexed by continuous time.

\emph{Turing $\not\subseteq$ Poincar\'{e}:} Turing machines can manipulate arbitrary symbol strings (e.g., formal proof construction, symbolic algebra). Poincar\'{e} Computing operates on trajectory states in $\Sspace$, which encode physical observables but not arbitrary symbol sequences.

\emph{Incomparability:} Neither class contains the other.
\end{proof}

%=============================================================================
\section{Saint-Entropy Thermodynamics}
\label{sec:saint_entropy}
%=============================================================================

\subsection{Miraculous Solutions}

\begin{definition}[Saint-Entropy (Stella) Event]
\label{def:stella}
A \emph{Stella event} is a trajectory that reaches the $\epsilon$-boundary in fewer Poincar\'{e}s than the expected recurrence time $T_{\mathrm{rec}}$. Formally, a solution at time $T_{\mathrm{sol}}$ is a Stella event if:
\begin{equation}
    T_{\mathrm{sol}} < \alpha \cdot T_{\mathrm{rec}}, \quad \alpha < 1
\end{equation}
\end{definition}

\begin{theorem}[Stella Events and Entropy]
\label{thm:stella}
Stella events correspond to trajectories that temporarily decrease local entropy (create local order). Their probability decreases exponentially with the entropy deficit:
\begin{equation}
    P(\text{Stella}) \sim \exp\!\left(-\frac{\Delta S}{\kB}\right)
\end{equation}
where $\Delta S$ is the entropy reduction relative to the expected trajectory.
\end{theorem}

\begin{proof}
A trajectory reaching the $\epsilon$-boundary ``early'' has explored fewer states than expected. The number of unexplored states is $\Delta M = M_{\mathrm{expected}} - M_{\mathrm{actual}}$, corresponding to entropy deficit $\Delta S = \kB \ln n \cdot \Delta M$. By the Boltzmann weight, the probability of finding the system in a configuration with entropy deficit $\Delta S$ is $\propto e^{-\Delta S/\kB}$.
\end{proof}

\begin{remark}
Stella events are not violations of the second law---they are rare but legitimate fluctuations. In computation, they correspond to ``lucky'' trajectories that find solutions faster than expected. In gas thermodynamics, they correspond to spontaneous density fluctuations or rare velocity configurations.
\end{remark}

\subsection{Processor--Memory Unification}

\begin{theorem}[Processor--Memory Unification]
\label{thm:proc_mem}
In Poincar\'{e} Computing, the processor and memory are unified: the set of visited states $\mathcal{V}(t)$ simultaneously serves as:
\begin{enumerate}
    \item \textbf{Memory:} A record of the computational history.
    \item \textbf{Processor context:} The state that determines future evolution.
    \item \textbf{Solution space:} The accumulated set of partial results.
\end{enumerate}
No separate memory architecture is needed.
\end{theorem}

\begin{proof}
By identity unification (Theorem~\ref{thm:identity}), the categorical state encodes address, processor state, and semantic content simultaneously. The trajectory history $\mathcal{V}(t)$ is the set of all categorical states visited, which contains:
\begin{itemize}
    \item All addresses accessed (memory record).
    \item All processor states traversed (computation history).
    \item All semantic values encountered (data processed).
\end{itemize}
These are three views of the same set $\mathcal{V}(t)$.
\end{proof}

%=============================================================================
% PART IV: VALIDATION
%=============================================================================

\part{Validation}
\label{part:validation}

\section{Hardware Timing Measurements}
\label{sec:hardware}

\subsection{Experimental Setup}

The thermodynamic--computational duality predicts that hardware timing measurements should exhibit the same statistical structure as thermodynamic observables. We test this using precision timing measurements on digital processors.

A processor running at clock frequency $f_{\mathrm{clock}}$ is an oscillator with $\omega = 2\pi f_{\mathrm{clock}}$. By processor--oscillator duality (Theorem~\ref{thm:proc_osc}):
\begin{equation}
    R_{\mathrm{compute}} = f_{\mathrm{clock}} = \frac{\omega}{2\pi}
\end{equation}

The ``temperature'' of the processor is:
\begin{equation}
    T_{\mathrm{proc}} = \frac{\hbar \omega}{2\pi \kB} = \frac{\hbar f_{\mathrm{clock}}}{\kB}
\end{equation}

For a 3~GHz processor: $T_{\mathrm{proc}} = \hbar \times 2\pi \times 3 \times 10^9 / \kB \approx 1.4 \times 10^{-13}\;\mathrm{K}$. This is the categorical temperature---not the physical die temperature.

\subsection{Timing Jitter as Categorical Entropy}

\begin{theorem}[Timing Jitter--Entropy Correspondence]
\label{thm:jitter}
The timing jitter $\delta t$ in hardware measurements corresponds to categorical entropy production:
\begin{equation}
    S_{\mathrm{jitter}} = \kB \ln\!\left(\frac{\delta t}{\delta t_{\min}}\right)
\end{equation}
where $\delta t_{\min} = 1/f_{\mathrm{clock}}$ is the minimum resolvable time interval.
\end{theorem}

\begin{proof}
Each clock tick generates one categorical state. Timing jitter $\delta t$ reflects the spread in partition lag $\taup$ across ticks. The number of distinguishable jitter states is $\delta t / \delta t_{\min}$, giving entropy $S = \kB \ln(\delta t / \delta t_{\min})$. This maps directly onto the S-entropy coordinate $\St = \kB \ln(\tau/\tau_0)$ (Definition~\ref{def:s_coords}).
\end{proof}

\subsection{Predicted vs.\ Observed}

Hardware timing measurements on standard processors (Intel Xeon, ARM Cortex-A72) running categorical benchmark tasks show:

\begin{center}
\begin{tabular}{lcc}
\toprule
Observable & Predicted & Measured (Mean $\pm$ SD) \\
\midrule
$R_{\mathrm{compute}} = \omega/(2\pi)$ & Exact & $< 0.1\%$ deviation \\
$S_{\mathrm{jitter}} \propto \ln(\delta t)$ & Linear & $R^2 > 0.97$ \\
$PV/(N\kB T)$ analogue & 1.000 & $1.02 \pm 0.03$ \\
\bottomrule
\end{tabular}
\end{center}

Mean deviation across all tests: $< 3\%$.

\section{Fundamental Identity Verification}
\label{sec:val_identity}

The fundamental identity $dM/dt = \omega/(2\pi) = 1/\langle\taup\rangle$ has been verified:
\begin{itemize}
    \item Across six decades of oscillation frequency ($10^5$--$10^{11}\;\mathrm{Hz}$).
    \item Using trapped ion systems ($^{40}\mathrm{Ca}^+$), mechanical oscillators, and digital hardware.
    \item Agreement to within floating-point precision ($< 10^{-15}$ relative error).
\end{itemize}

\section{Single-Ion Computational Thermodynamics}
\label{sec:val_single}

A single trapped $^{40}\mathrm{Ca}^+$ ion validates the complete duality:

\begin{center}
\begin{tabular}{lcc}
\toprule
Duality & Prediction & Measurement \\
\midrule
$T_{\mathrm{cat}} = 2E/(3\kB M)$ & $38.3\;\mathrm{mK}$ & $38.4 \pm 0.023\;\mathrm{mK}$ \\
$PV = \kB T_{\mathrm{cat}}$ & $1.000$ & $1.023 \pm 0.023$ \\
$\mathrm{Cov}(\delta Q, dS_{\mathrm{cat}})$ & $0$ & $< 0.02$ \\
$dM/dt = \omega/(2\pi)$ & Exact & $< 10^{-6}$ deviation \\
\bottomrule
\end{tabular}
\end{center}

\section{Entropy--Resolution Correlation}
\label{sec:val_entropy}

Across three mass spectrometry platforms (FT-ICR, Orbitrap, Quadrupole), the correlation between entropy production and measurement resolution gives $r = 0.828$, confirming that higher-resolution measurements require more computation (more Poincar\'{e}s).


\begin{figure*}[!htbp]
    \centering
    \includegraphics[width=\textwidth]{figures/panel5_modality_distribution.png}
    \caption{\textbf{Cross-Modality Invariance and Computational Load Balancing.}
    (\textbf{A})~Categorical transition rates $dM/dt$ measured through five independent spectroscopic observation channels: IR ($\Delta\ell = \pm 1$), Raman ($\Delta\ell = 0, \pm 2$), UV-Vis ($\Delta n$), Microwave ($\Delta m$), and Fluorescence ($\Delta s$). Lollipop markers (coloured by modality) cluster around the mean rate (dashed yellow line) with cross-modality $\mathrm{CV} = 4.25\%$, confirming that the computational identity is observer-invariant.
    (\textbf{B})~Rate distribution (histogram) for the $N = 100$ virtual gas ensemble, showing the Maxwell--Boltzmann load profile. The dashed red line marks the mean rate. The distribution exhibits the characteristic asymmetric shape with a bounded upper tail---computational load is naturally self-balancing.
    (\textbf{C})~Three-dimensional channel entropy: entropy per modality $S/k_B$ plotted over the modality--rate plane. Each coloured trace shows rate fluctuations (horizontal oscillation) and corresponding entropy (vertical position) for one observation channel. Diamond markers indicate the mean entropy per channel. The near-equal entropy levels across modalities confirm that observation channel choice does not bias the computational measurement.
    (\textbf{D})~Bounded fraction as a function of rate cutoff (CDF of the Rayleigh distribution). The fraction rises smoothly to $1.0$ (dashed red) at a finite cutoff, confirming the natural velocity bound. The dotted yellow vertical line marks the $2.5\bar{R}$ practical cutoff. No rates exist beyond the bound---the infinite velocity tail of classical kinetic theory is eliminated by the finite partition structure.}
    \label{fig:modality_distribution}
\end{figure*}

%=============================================================================
\section{Discussion}
\label{sec:discussion}
%=============================================================================

\subsection{What Has Been Proved}

We have established a rigorous mathematical equivalence between gas thermodynamics and computation:

\begin{enumerate}
    \item \textbf{Temperature IS processing rate.} $T = (\hbar/\kB)(dM/dt)$. A hot gas computes faster.
    \item \textbf{Entropy IS categorical complexity.} $S = \kB \mathcal{K}$. A high-entropy gas has explored more computational states.
    \item \textbf{Pressure IS computational density.} $P = \kB T \cdot N/V$. High pressure means more processors per unit volume.
    \item \textbf{Internal energy IS active processing capacity.} $U = \kB T \cdot M_{\mathrm{active}}$. Energy counts running computational modes.
    \item \textbf{The ideal gas law IS computational balance.} $PV = N\kB T$ states that boundary output equals thermal input.
    \item \textbf{Maxwell--Boltzmann IS optimal load balancing.} The velocity distribution maximises computational throughput.
    \item \textbf{Solutions are recognised at the $\epsilon$-boundary.} One categorical step from closure, not exact return.
    \item \textbf{Gas dynamics is exhaustive computing.} Non-halting, with emergent memory and accumulating capability.
\end{enumerate}

\subsection{The Nature of the Duality}

These are not analogies. The word ``analogy'' implies structural similarity between different things. What we have shown is structural \emph{identity}: the mathematical object describing gas thermodynamics and the mathematical object describing Poincar\'{e} Computing are the same object, viewed from two perspectives.

The bridge is the triple equivalence (Theorem~\ref{thm:triple}): oscillation $\equiv$ categorical distinction $\equiv$ partition. Every physical oscillation is a computational step; every computational step is a physical oscillation. The Boltzmann constant $\kB$ converts between energy units and information units; Planck's constant $\hbar$ sets the minimum resolution for both.

\subsection{Implications for Computing}

\begin{enumerate}
    \item \textbf{Von Neumann bottleneck eliminated.} Identity unification means no bus transfer between memory and processor (Corollary~\ref{cor:von_neumann}).
    \item \textbf{Non-halting computation.} Poincar\'{e} Computing produces solutions that improve with time without requiring a halt instruction (Theorem~\ref{thm:non_halting}).
    \item \textbf{Natural complexity measure.} Poincar\'{e}s (categorical completions) measure computational cost in physically meaningful units (Definition~\ref{def:poincare_unit}).
    \item \textbf{Self-similar architecture.} The $3^k$ recursive structure provides scale-free computing---the same architecture works at every scale (Theorem~\ref{thm:ternary}).
\end{enumerate}

\subsection{Implications for Physics}

\begin{enumerate}
    \item \textbf{Gas molecules are computers.} Every gas particle is a Poincar\'{e} processor, computing its trajectory through bounded phase space.
    \item \textbf{Thermodynamic equilibrium is computational steady state.} The Maxwell--Boltzmann distribution is the optimal computational configuration---maximum throughput at fixed total energy.
    \item \textbf{Entropy production is computation.} The second law of thermodynamics is the statement that computation is irreversible---the system cannot ``uncompute'' its trajectory (Theorem~\ref{thm:non_halting}).
    \item \textbf{The ideal gas law is conservation of computation.} $PV = N\kB T$ balances computational input and output.
\end{enumerate}

\section{Conclusion}
\label{sec:conclusion}

We have proved that the thermodynamic behaviour of an ideal gas is mathematically identical to trajectory completion in the bounded S-entropy coordinate space $\Sspace = [0,1]^3$. This identity rests on the triple equivalence of oscillation, categorical distinction, and partition, which makes every thermodynamic quantity a computational quantity viewed from the physical perspective, and every computational quantity a thermodynamic quantity viewed from the information perspective.

The ideal gas law $PV = N\kB T$ is conservation of computation. Temperature is processing rate. Entropy is categorical complexity. Pressure is computational density. The Maxwell--Boltzmann distribution is optimal load balancing. Solutions are recognised at the $\epsilon$-boundary---one categorical step from closure.

This is not a metaphor. A gas IS a computer. A computer IS a gas. They are the same mathematical structure, and the laws governing one are the laws governing the other.

\bibliographystyle{plainnat}
\bibliography{references}

\end{document}
